\Extrachap{Next Steps}

There are several next steps to this book. This is not even a first edition. This is a pre-edition, a third draft of what will be a completed book. One thing is clear though - it will always be a self-published, no cost (eBook) or low cost (printing cost only) textbook. Though as the pages increase the cost for pages will increase. But no fear. If you've purchased this draft, keep it and if a new draft or when the first edition is printed, I'll give you a new printed copy \textit{gratis}. 

\subsection*{New Content}

\textbf{Chapter 13 - Digestion, Absorption and Metabolism} is currently a missing but necessary chapter for \textbf{Part III} on Muscle Support. Without food and water entering the body, being digested, absorbed and metabolized the extra-cellular fluid cannot sustain muscle fibers. It has not been written yet because - at least at Plymouth State University DPT - it is not part of the first Clinical Physiology course. It is taught in the Nutrition and Exercise Prescription course. It is forthcoming and in progress. It sits at the border of Part III on muscle support and Part IV on the maintenance of long term muscle cell health, growth and repair.

\textbf{Part IV} is currently in process and focuses on muscle molecular biology. Specifically, how the muscle continues to fulfill its role of creating tension without being too large of a burden on the physiological support systems. Topics include maintaining muscle cells and tissue for the demands they are placed under through the processes of atrophy, hypertrophy, and repair. 

\textbf{Part V} is also in process and will include a variety of integrated exercises. These topics have been selected to highlight the need for physiological coordination such as Fick's Equation to highlight integration of the material in Parts II and III; exercise prescription and antigravity, cold and hypoxia exposure to highlight integration of the material in Parts II, III and IV; the autonomic nervous system to highlight integration of Parts III and IV; and rehabilitation protocols to highlight integration of Part I with the all other parts as the full spectrum of hierarchy is manifest in clinical practice.

Clinical practice is a fulfillment of clinical physiology. Practice abductively reasons based on physiology to understand movement problems; and practice attempts to deductively modify physiology through interventions.

\subsection*{Update Content}

It is hoped that all of the typos, errors and mishaps in the existing content will be fixed prior to the final publication of the first edition of the text. There are still some rather long and muddy sections that require at least one more pass through in the tried and tested locale of classroom teaching and student feedback!

\subsection*{Added Features}

Additional features for the first edition will be the inclusion of a glossary of key terms at the end of each chapter, and an index at the end of the book. In all honesty, I really need a full printed draft to highlight those terms worty of being in the glossary and the index.
Finally, an end of book Works Cited page is necessary for those - like me - that enjoy reading through the works cited in masse and not just by chapter. This feature will not remove the end of chapter references because both are options worthy of inclusion for different preferences.
